\markboth{INTRODUCTION}{INTRODUCTION}
\section{Introduction}\label{sec:intro}

Wearable devices from Garmin, Polar, and Apple continuously capture heart rate, GPS tracks, sleep stages, heart-rate variability (HRV), and stress~\citep{seckin_review_2023}, and Strava alone holds activity data from over 180~million users~\citep{strava_trends_2025}. Yet all this training data stays split across proprietary ecosystems with incompatible interfaces, so analysing it as a whole is hard for anyone without sports-science training~\citep{bourdon_monitoring_2017}. An endurance cyclist we interviewed shows the problem: after every ride he exports his Strava data and uploads it to ChatGPT, which is slow and incomplete, as the general-purpose model cannot see his Garmin recovery data, the weather forecast, or his calendar. Deciding whether to train on a given day instead means checking a Garmin watch for HRV and sleep, Strava for cumulative load, a weather service for conditions, a calendar for scheduling, and stitching the four together is left entirely to the athlete.

Large language models (LLM) can reason across different data sources once they are able to call tools~\citep{schick_toolformer_2023, qu_tool_2024}. Two new open protocols now make it feasible to connect these fragmented sources behind one conversational interface. Anthropic's Model Context Protocol (MCP)~\citep{anthropic_mcp_2024} for tool discovery and invocation, and Google's Agent-to-Agent protocol (A2A)~\citep{google_a2a_2025} for inter-agent coordination.

\subsection{Problem Statement and Approach}\label{sec:intro:problem}

This paper presents \textit{Training Copilot}, an open-source sports analytics platform addressing three shortcomings of the current landscape. \textit{Data fragmentation}: training and wellness data is spread across platforms with incompatible authentication, formats, and access patterns, and although sports-science consensus holds that effective load monitoring must combine external and internal load from multiple sources~\citep{bourdon_monitoring_2017, halson_monitoring_2014}, the tools that collect it provide no unified view. Training Copilot solves this with MCP as a uniform data layer: each external API is an independent MCP server, and a single host client discovers and routes tool calls without provider-specific logic. A \textit{Lack of cross-domain reasoning} is evident in platforms such as Strava, Garmin Connect, Intervals.icu, and TrainingPeaks display metrics in isolation, even though readiness to train depends on the interplay of recovery, workload, weather, and schedule~\citep{bourdon_monitoring_2017, gabbett_training_2016}. Training Copilot addresses this with six LangGraph~\citep{langchain_langgraph_2024} ReAct~\citep{yao_react_2023} agents, each scoped to a domain (recovery, load, context, routes, fitness knowledge, goal-driven coaching) and coordinated by an orchestrator via A2A. Existing platforms have \textit{Rigid architectures} and are tightly coupled to their data sources. Adding a new source, such as a nutrition tracker, therefore requires code changes throughout the stack. Training Copilot's design reduces this to one server file and one configuration line, consistent with microservice principles~\citep{dragoni_microservices_2017}.  

\subsection{Contributions}\label{sec:intro:contributions}

\begin{itemize}
\item An MCP-based integration architecture achieving single-line extensibility, verified empirically by adding a Telegram bridge and a flythrough renderer without changes to the host, agents, or User Interface (UI).
\item A multi-agent system with six domain-scoped specialists coordinated via A2A; scoping each specialist to at most three MCP servers counters the tool-selection degradation observed with large tool sets~\citep{qu_tool_2024}.
\item A RAG-based fitness knowledge agent that grounds exercise-science advice in a local vector database of sports-science literature, with explicit citation preservation.
\item A goal-driven coaching layer: a structured race goal with milestones, deterministically computed training zones, and a guardrail-validated multi-week training plan, in which all training arithmetic is plain, user-checkable code and the LLM only selects and explains workouts.
\item An automated evaluation framework built on MLflow's multi-turn evaluation tooling~\citep{mlflow_multiturn_2025} that scores the logged conversations of real users on five dimensions, including a deterministic trace-based grounding scorer.
\end{itemize}

\subsection{Structure of This Paper}\label{sec:intro:structure}

Section~\ref{sec:market} analyses the sports analytics market and the gap Training Copilot addresses; Section~\ref{sec:sota} reviews the state of the art in agentic AI, tool-integration protocols, sports analytics, and retrieval-augmented generation; Section~\ref{sec:architecture} presents the system architecture; Section~\ref{sec:data} describes the data sources; Section~\ref{sec:agents} details agent interaction patterns; Section~\ref{sec:evaluation} describes the evaluation methodology; Section~\ref{sec:discussion} discusses trade-offs, limitations, and ethics; and Section~\ref{sec:conclusion} closes with a summary and future work.
